ChunkLoris sits at the intersection of two long lineages: low-bandwidth
application-layer denial of service (the slow-DoS family) and
algorithmic-complexity / resource-asymmetry attacks. Our contribution is
not a new primitive but a \emph{measurement and systematization}: we
identify and quantify a resource axis---per-DATA-unit application
\emph{dispatch} CPU cost---that the prior literature names only in
passing, and we measure it across \num{27} production servers. We
organize the prior work by the resource each attack amplifies, and draw
the line to ours in each case.

\subsection{Connection-slot exhaustion: the slow-DoS family}
\label{sec:rel-slowdos}

The canonical low-and-slow attacks amplify \emph{connection slots}: an
attacker holds open many concurrent connections at negligible client
cost, exhausting a server's bounded pool. Slowloris~\cite{slowloris}
(Hansen, 2009) keeps thread-per-connection servers' worker pools occupied
by dribbling partial request headers and never completing them; RUDY
(R-U-Dead-Yet)~\cite{rudy} does the same with a slow \texttt{POST} body
under a large \fname{Content-Length}; slow-read variants stall on the
response side. Cambiaso et al.~\cite{cambiaso_survey} give the standard
taxonomy of these attacks (pending-request, long-response, and
multi-layer slow DoS). Section~\ref{sec:bg-cwe} situates ChunkLoris in
this lineage at the wire level.

The distinction is the bottleneck. The slow-DoS family is
\emph{concurrency-bound}: a coalescing front-end or a per-connection
header/idle timeout neutralizes it, and once admitted the request costs
the server essentially nothing. ChunkLoris is \emph{CPU-bound}: a single
admitted, RFC-compliant, fully-framed chunked request burns single-core
CPU \emph{in proportion to its chunk count} on the parser/dispatcher
boundary, independent of how many connections the attacker holds. Header
timeouts and slot caps do not bound it; only per-request work or byte
accounting does.

\subsection{Algorithmic-complexity and asymmetric-resource attacks}
\label{sec:rel-complexity}

Crosby and Wallach~\cite{crosby2003} introduced low-bandwidth
algorithmic-complexity attacks: worst-case inputs that drive a data
structure (hash tables, binary trees) into its degenerate
$\Theta(n^2)$ regime, converting a small request into large server CPU.
This is the conceptual ancestor of the superlinear case we measure
(Node's \fwid{http.Server} via \fname{process.nextTick}). MITRE
catalogues this class as CWE-407 (Inefficient Algorithmic
Complexity)~\cite{cwe407}, the asymmetry itself as CWE-405 (Asymmetric
Resource Consumption (Amplification))~\cite{cwe405}, and the umbrella as
CWE-400 (Uncontrolled Resource Consumption)~\cite{cwe400}.

ChunkLoris's \emph{common} (linear) case is not an algorithmic-complexity
bug in the Crosby--Wallach sense: the per-chunk cost is $\Theta(N)$ in
chunk count with a large constant, not a degenerate worst case of an
otherwise sub-quadratic structure. The amplified resource is the
\emph{fixed per-dispatch overhead}---a parser callback plus an
application-receive event (and, on most event loops, a task/await
switch)---multiplied by an attacker-chosen frame count. It is best read
as CWE-405 asymmetry whose unit of amplification is the protocol's
smallest framing unit rather than a hash collision. The three Node
servers additionally exhibit CWE-407 behavior, which we treat as a
distinct, server-specific finding layered on top of the universal linear
tax.

\subsection{Control-plane and decompression attacks in HTTP/2}
\label{sec:rel-h2}

The modern HTTP/2 DoS literature amplifies resources other than
per-DATA-unit dispatch, which sharpens our novelty claim. \emph{Rapid
Reset} (CVE-2023-44487)~\cite{rapidreset2023,rapidreset_cloudflare}
amplifies \emph{control-plane stream churn}: a client opens and
immediately \fname{RST\_STREAM}s streams, decoupling expensive
server-side request setup from the 100-concurrent-stream cap and
producing record-breaking request rates. The \emph{CONTINUATION Flood}
(CVE-2024-27316 and siblings)~\cite{continuation2024,continuation_nowotarski}
amplifies \emph{memory and header-decode CPU} via an unbounded stream of
\fname{CONTINUATION} frames that are decoded but never finalized.
\emph{MadeYouReset} (CVE-2025-8671)~\cite{madeyoureset,madeyoureset_cert}
revives the Rapid Reset control-plane axis by inducing
\emph{server}-initiated resets through malformed control frames, thereby
sidestepping client-reset rate limits. The calif.io ``hidden HTTP/2
bomb''~\cite{calif_http2bomb} composes an HPACK indexed-reference bomb
(one wire byte per reference expands to tens-to-thousands of bytes of
server allocation) with a flow-control window stall, amplifying
\emph{memory via header decompression} and pinning it with zero-window
\fname{WINDOW\_UPDATE} drips.

None of these targets per-DATA-frame application dispatch. Rapid Reset and
MadeYouReset never send \fname{DATA}; CONTINUATION Flood and the HPACK bomb
live entirely in the header/control plane. Crucially, HTTP/2 flow control
bounds \emph{bytes} in flight but places no ceiling on the \emph{number of
frames} carrying those bytes---so an attacker can split one
flow-control-window's worth of body into thousands of \SI{1}{\byte}
\fname{DATA} frames, each forcing a parser callback and a handler wake. Our
Wave 3 results show every HTTP/2 server in our sample inheriting the same
per-frame tax (\SIrange{4}{10}{\micro\second}/frame) as HTTP/1.1
chunked-TE; the per-chunk axis is orthogonal to, and unmitigated by, the
defenses deployed for the attacks above.

\subsection{HTTP parsing and smuggling measurement lineage}
\label{sec:rel-smuggling}

Methodologically, our cross-server probing of how independent
implementations diverge on an underspecified part of the HTTP grammar
follows the request-smuggling measurement tradition, most prominently
Kettle's \emph{HTTP Desync Attacks}~\cite{kettle_desync}, which exploits
\emph{disagreement} between front-end and back-end parsers over message
framing (\fname{Content-Length} vs.\ \fname{Transfer-Encoding}). That line
weaponizes \emph{semantic} parser divergence to smuggle requests; we
instead measure \emph{performance} divergence on a point the spec leaves
open---the granularity at which a conformant parser delivers decoded body
bytes to the handler (Section~\ref{sec:bg-chunked}). We share the
comparative, many-implementation methodology but target a CPU cost axis
rather than a confusion-based exploit, and every server we measure is
fully RFC-compliant.
